% Answers to Chapter 4, from lectures/L04/appendix/c_solutions.md; the self-test from
% lectures/L04/README.md.

\facitavsnitt{c4}{Kapitel 4}{Ekvationssystem}
\begin{facit}

\svar{1.1} \sv{a} $9{,}4\,\text{V}$ \sv{b} $1{,}88\,\text{W}$

\svar{1.2} \sv{a} $R = \frac{\tau}{C} = 2\,000\,\Omega$
\sv{b} $2\,\text{k}\Omega$

\svar{1.3} $5U_{\text{in}} \leq 15$, alltså $U_{\text{in}} \leq 3\,\text{V}$.

\svar{2.1} $x = \frac{20}{7} \approx 2{,}86$, $y = \frac{17}{7} \approx 2{,}43$

\svar{2.2} \sv{a} $I_1 + I_2 = 5$ och $I_1 = 2I_2$
\sv{b} $I_1 = \frac{10}{3} \approx 3{,}33\,\text{A}$, $I_2 = \frac{5}{3} \approx 1{,}67\,\text{A}$

\svar{2.3} \sv{a} $U_1 + U_2 = 12$ och $U_2 = 3U_1$
\sv{b} $U_1 = 3\,\text{V}$, $U_2 = 9\,\text{V}$
\sv{c} $R_1 = 1{,}5\,\Omega$, $R_2 = 4{,}5\,\Omega$

\svar{3.1} \sv{a} $(2, 5)$ \sv{b} $(9, 3)$ \sv{c} $(7, 3)$

\svar{3.2} \sv{a} $(3, 1)$ \sv{b} $(4, 2)$ \sv{c} $(3, 2)$ \sv{d} $(2, 1)$

\svar{3.3} \sv{a} Oändligt många: ekvation (2) är ekvation (1) gånger $2$, så linjerna
sammanfaller.
\sv{b} Ingen: samma vänsterled men olika högerled, så linjerna är parallella.
\sv{c} Exakt en, $(2, 2)$: linjerna har olika lutning och skär varandra.
\sv{d} Oändligt många: ekvation (2) gånger $-2$ är ekvation (1), så linjerna sammanfaller.

\svar{3.4} \sv{a} För $x = 0, 1, 2, 3, 4$ ger $y = 2x - 1$ värdena $-1, 1, 3, 5, 7$ och
$y = -x + 5$ värdena $5, 4, 3, 2, 1$.
\sv{b} Vid $x = 2$; skärningspunkten är $(2, 3)$.
\sv{c} $(x, y) = (2, 3)$, samma som avläsningen.
\sv{d} Raderna skulle aldrig ge samma $y$: skillnaden mellan dem vore konstant, som för två
parallella linjer.

\svar{3.5} \sv{a} $(4, 2)$ \sv{b} $(6, 3)$

\svar{3.6} \sv{a} $I_1 + I_2 = 12$ och $I_1 - I_2 = 4$
\sv{b} $I_1 = 8\,\text{A}$, $I_2 = 4\,\text{A}$
\sv{c} $8 + 4 = 12$ och $8 - 4 = 4$ stämmer.

\svar{3.7} \sv{a} $I_1 + I_2 = 3$ och $I_1 = 2I_2$
\sv{b} $I_1 = 2\,\text{A}$, $I_2 = 1\,\text{A}$
\sv{c} $R_1 = 12\,\Omega$, $R_2 = 24\,\Omega$
\sv{d} $R_1 \parallell R_2 = \frac{12 \times 24}{12 + 24} = 8\,\Omega = \frac{24}{3}\,\Omega$

\svar{3.8} \sv{a} $U_1 + U_2 = 15$ och $U_1 - U_2 = 3$
\sv{b} $U_1 = 9\,\text{V}$, $U_2 = 6\,\text{V}$
\sv{c} $R_1 = 30\,\Omega$, $R_2 = 20\,\Omega$
\sv{d} $R_1 + R_2 = 50\,\Omega = \frac{15}{0{,}3}\,\Omega$

\svar{3.9} \sv{a} $U = 12 - 2I$ och $U = 4I$
\sv{b} $I = 2\,\text{A}$, $U = 8\,\text{V}$
\sv{c} $12 - 2 \times 2 = 8$ och $4 \times 2 = 8$ stämmer.
\sv{d} Arbetspunkten.

\svar{3.10} \sv{a} $\frac{6}{9} = \frac{R_1}{900}$, som ger $R_1 = 600\,\Omega$.
\sv{b} $R_1 = 600\,\Omega$, $R_2 = 300\,\Omega$
\sv{c} $U_1 = 9 \times \frac{600}{900} = 6\,\text{V}$ stämmer.

\svar{3.11} \sv{a} $I_1 = 3\,\text{A}$, $I_2 = 2\,\text{A}$, $I_3 = 4\,\text{A}$
\sv{b} $3 + 2 + 4 = 9$, $3 - 2 = 1$ och $I_3 = 4$ stämmer.
\sv{c} Den ger en obekant direkt; insatt i de andra blir det ett system med två obekanta.

\svar{3.12} \sv{a} $5R_A + 3R_B = 4\,390$ och $2R_A + 4R_B = 3\,660$
\sv{b} $R_A = 470\,\Omega$, $R_B = 680\,\Omega$
\sv{c} $2\,350 + 2\,040 = 4\,390$ och $940 + 2\,720 = 3\,660$ stämmer.

\svar{3.13} \sv{a} Ja. \sv{b} Ja. \sv{c} Nej: $5 - 2 \times 1 = 3$ stämmer, men $5 + 1 = 6 \neq 7$.
\sv{d} Lösningen måste ligga på båda linjerna samtidigt. Ett talpar som bara uppfyller den ena
ekvationen ligger bara på den ena linjen, som i c).

\svar{3.14} \sv{a} $(2, 3)$ \sv{b} $(4, -1)$ \sv{c} $(2, -1)$
\sv{d} Substitution i a) och b), där en obekant enkelt kan lösas ut; addition i c), där
$y$-termerna enkelt blir motsatta. Båda metoderna ger samma svar.

\svar[Testa dig själv]{T} \sv{1} $(x, y) = (3, 2)$
\sv{2} $I_1 = 3\,\text{A}$, $I_2 = 1\,\text{A}$

\end{facit}
